% ----------------------------------------------------------------------
% Capítulo 1: Introdução
% ----------------------------------------------------------------------
\chapter{Introdução}
\label{cap:introducao}

\instrucao{De três a cinco páginas. Apresente o problema prático que motivou o trabalho, o contexto profissional em que ele aparece, os objetivos e a estrutura da dissertação. No mestrado profissional, deixe claro desde o início qual produto técnico será entregue.}

As bibliotecas escolares da rede pública recebem livros de várias fontes ao mesmo tempo: programas federais, doações da comunidade e compras pontuais da escola. Sem uma política escrita, a seleção depende de quem está no balcão naquele ano, e o descarte quase nunca acontece. A Lei nº 12.244, de 24 de maio de 2010, ao tratar da universalização das bibliotecas escolares, reforçou a necessidade de acervos organizados e atualizados \cite{brasil2010}.

\section{Problema}
\instrucao{Escreva o problema em uma frase interrogativa, ligada à prática. Exemplo: como sistematizar critérios de seleção e descarte que equipes pequenas consigam aplicar?}

Como sistematizar critérios de seleção, aquisição, avaliação e descarte que as equipes das bibliotecas escolares municipais consigam aplicar sem apoio externo permanente?

\section{Objetivos}
\subsection{Objetivo geral}
Elaborar e validar um guia de política de desenvolvimento de coleções para bibliotecas escolares da rede pública municipal.

\subsection{Objetivos específicos}
\instrucao{Use verbos no infinitivo que possam ser verificados ao final: diagnosticar, elaborar, validar, propor.}
\begin{enumerate}[label=\alph*)]
  \item diagnosticar as práticas de formação de acervo das bibliotecas participantes;
  \item identificar, na literatura, os elementos mínimos de uma política de desenvolvimento de coleções;
  \item elaborar o guia como produto técnico;
  \item validar o guia com as equipes, em oficinas.
\end{enumerate}

\section{Justificativa}
\instrucao{Mostre a relevância para a instituição e para a área. No mestrado profissional, o impacto esperado do produto conta tanto quanto a lacuna teórica.}

A política escrita reduz decisões improvisadas, dá transparência ao uso de recursos públicos e protege a biblioteca diante de doações inadequadas. O guia proposto pode ser adotado por outras redes com ajustes mínimos.

\section{Estrutura do trabalho}
O Capítulo~\ref{cap:referencial} reúne os fundamentos. O Capítulo~\ref{cap:metodologia} descreve o percurso metodológico. O Capítulo~\ref{cap:diagnostico} apresenta o diagnóstico e a validação do produto. O Capítulo~\ref{cap:consideracoes} traz as considerações finais. O guia completo está no Apêndice~\ref{ap:produto}. A apresentação segue a NBR 14724 \cite{NBR14724:2011}, e as referências seguem a NBR 6023 \cite{NBR6023:2018}.
